\section{Timeline IR}
\label{sec:ir}

The Timeline Intermediate Representation (IR) is the central contract of this system. It defines
\emph{what} a timeline communicates---entities, relationships, temporal structure, and semantic
status---while remaining agnostic to \emph{how} that communication is rendered. This separation
enables deterministic rendering, source-control-friendly editing, and reliable generation by
both humans and AI agents.

\subsection{Design Principles}

\begin{enumerate}
  \item \textbf{Make illegal states unrepresentable.} The IR's type system should prevent
        malformed timelines at the structural level, not merely validate them after the fact.
  \item \textbf{Rendering-agnostic.} The IR describes meaning, not pixels. Visual
        decisions---colors, fonts, positioning---belong to the theme engine.
  \item \textbf{Git-friendly.} Stable field ordering, line-oriented YAML syntax, deterministic
        serialization, minimal diff churn.
  \item \textbf{Agent-generatable.} Clear required-vs-optional distinctions, explicit defaults,
        forgiving structure that remains validatable.
  \item \textbf{Orthogonal core.} A small set of primitives; convenience features are optional
        layers, not core complexity.
\end{enumerate}

\subsection{Type Vocabulary}
\label{sec:ir-types}

The IR uses the following type vocabulary consistently throughout this specification:

\begin{table}[h]
\centering
\begin{tabular}{lp{10cm}}
\toprule
\textbf{Type} & \textbf{Description} \\
\midrule
\texttt{string} & UTF-8 text, non-empty unless explicitly noted \\
\texttt{string?} & Optional string (may be omitted or null) \\
\texttt{int} & Integer \\
\texttt{number} & Decimal number \\
\texttt{number[0..1]} & Decimal in closed interval $[0, 1]$ \\
\texttt{bool} & Boolean (\texttt{true}/\texttt{false}) \\
\texttt{date} & A temporal point (see \S\ref{sec:date-model}) \\
\texttt{daterange} & A temporal interval with start and optional end \\
\texttt{id} & Document-unique identifier (kebab-case slug, e.g., \texttt{alpha-release}) \\
\texttt{ref<T>} & Reference to an entity of type \texttt{T} by its \texttt{id} \\
\texttt{list<T>} & Ordered list of elements of type \texttt{T} \\
\texttt{optional<T>} & Value of type \texttt{T} that may be omitted \\
\texttt{enum\{...\}} & One of the listed literal values \\
\texttt{map<K,V>} & Key-value mapping \\
\bottomrule
\end{tabular}
\caption{IR type vocabulary}
\label{tab:type-vocab}
\end{table}

\subsection{Document Structure}
\label{sec:ir-structure}

A Timeline IR document is a single YAML file with the following root structure:

\begin{lstlisting}[language=yaml,caption={Root document structure}]
version: "1.0"
metadata: { ... }
tracks: [ ... ]
groups: [ ... ]       # optional
activities: [ ... ]
milestones: [ ... ]   # optional
annotations: [ ... ]  # optional
sections: [ ... ]     # optional
legend: { ... }       # optional
\end{lstlisting}

\subsubsection{Root Fields}

\begin{table}[h]
\centering
\small
\begin{tabular}{llccp{5cm}}
\toprule
\textbf{Field} & \textbf{Type} & \textbf{Req} & \textbf{Default} & \textbf{Description} \\
\midrule
\texttt{version} & \texttt{string} & Yes & --- & Specification version (semver, e.g., \texttt{"1.0"}) \\
\texttt{metadata} & \texttt{Metadata} & Yes & --- & Document-level metadata \\
\texttt{tracks} & \texttt{list<Track>} & Yes & --- & Swimlanes/rows; at least one required \\
\texttt{groups} & \texttt{list<Group>} & No & \texttt{[]} & Logical groupings of activities \\
\texttt{activities} & \texttt{list<Activity>} & Yes & --- & Time-spanning items; may be empty \\
\texttt{milestones} & \texttt{list<Milestone>} & No & \texttt{[]} & Point-in-time markers \\
\texttt{annotations} & \texttt{list<Annotation>} & No & \texttt{[]} & Callouts, notes, highlights \\
\texttt{sections} & \texttt{list<Section>} & No & \texttt{[]} & Visual/logical document sections \\
\texttt{legend} & \texttt{Legend} & No & auto & Legend mapping; auto-generated if omitted \\
\bottomrule
\end{tabular}
\caption{Root document fields}
\label{tab:root-fields}
\end{table}

\textbf{Invariant:} The document must contain at least one track. Activities and milestones
may both be empty, but a timeline with neither is semantically vacuous (valid but degenerate).


%% ============================================================================
%% METADATA
%% ============================================================================
\subsection{Metadata}
\label{sec:ir-metadata}

The \texttt{metadata} object contains document-level information and temporal framing.

\begin{table}[h]
\centering
\small
\begin{tabular}{llccp{4.5cm}}
\toprule
\textbf{Field} & \textbf{Type} & \textbf{Req} & \textbf{Default} & \textbf{Description} \\
\midrule
\texttt{title} & \texttt{string} & Yes & --- & Primary document title \\
\texttt{subtitle} & \texttt{string?} & No & \texttt{null} & Secondary title or tagline \\
\texttt{author} & \texttt{string?} & No & \texttt{null} & Author or team name \\
\texttt{created} & \texttt{date} & No & now & Document creation date \\
\texttt{updated} & \texttt{date} & No & now & Last modification date \\
\texttt{time\_range} & \texttt{TimeRange} & Yes & --- & Temporal bounds of the timeline \\
\texttt{axis\_unit} & \texttt{AxisUnit} & No & inferred & Granularity of the time axis \\
\texttt{theme} & \texttt{string?} & No & \texttt{null} & Theme identifier (resolved by renderer) \\
\texttt{locale} & \texttt{string?} & No & \texttt{"en"} & BCP-47 locale for date formatting \\
\texttt{today} & \texttt{date?} & No & eval time & Reference date for \texttt{now}, relative dates, and today-marker; see \S\ref{sec:date-anchor} \\
\texttt{fiscal\_year\_start} & \texttt{int [1..12]} & No & \texttt{1} & Calendar month on which the fiscal year begins (1~=~January); see \S\ref{sec:date-anchor} \\
\texttt{description} & \texttt{string?} & No & \texttt{null} & Extended description or notes \\
\bottomrule
\end{tabular}
\caption{Metadata fields}
\label{tab:metadata-fields}
\end{table}

\subsubsection{TimeRange}

The \texttt{time\_range} defines the temporal viewport of the timeline---the span that will
be rendered. Activities or milestones outside this range are valid but may be clipped.

\begin{lstlisting}[language=yaml]
time_range:
  start: 2026-Q1
  end: 2026-Q4
\end{lstlisting}

\begin{table}[h]
\centering
\begin{tabular}{llccp{5cm}}
\toprule
\textbf{Field} & \textbf{Type} & \textbf{Req} & \textbf{Default} & \textbf{Description} \\
\midrule
\texttt{start} & \texttt{date} & Yes & --- & Inclusive start of the viewport \\
\texttt{end} & \texttt{date} & No & open & Inclusive end; omit for open-ended \\
\bottomrule
\end{tabular}
\caption{TimeRange fields}
\label{tab:timerange-fields}
\end{table}

\subsubsection{AxisUnit}

\begin{lstlisting}
enum AxisUnit { day, week, month, quarter, half, year }
\end{lstlisting}

If omitted, the renderer infers axis unit from the time range span. The axis unit affects
visual tick marks and labels, not the precision of individual dates.

\subsubsection{Date Anchor and Fiscal Calendar}
\label{sec:date-anchor}

\paragraph{Date anchor (\texttt{today}).}
The symbolic date \texttt{now} and all relative date offsets (\texttt{+3m}, \texttt{-2w},
etc.) resolve against the \emph{date anchor}, evaluated in this order:
\begin{enumerate}
  \item \texttt{metadata.today} — if present, this value is used exclusively.
  \item \texttt{metadata.created} — fallback when \texttt{today} is absent.
  \item \textbf{Hard error} — if neither is present and any \texttt{now} or relative date
        appears in the document, the document is not deterministically evaluable and
        validation must reject it.
\end{enumerate}

Renderers \textbf{must not} consult the system clock to resolve \texttt{now} or relative
dates. The today-marker annotation (\texttt{type: today-marker}) also uses this anchor.

\textbf{Determinism caveat:} When \texttt{today} is omitted and \texttt{created} serves
as the anchor, output is deterministic (the \texttt{created} date is fixed in the
document). However, the document's visual "current position" does not advance with real
time. Authors and agents \textbf{should} set \texttt{today} explicitly whenever \texttt{now}
or relative dates appear, so that the document's meaning is unambiguous and byte-stable
across environments and rendering runs.

\paragraph{Fiscal calendar (\texttt{fiscal\_year\_start}).}
\texttt{fiscal\_year\_start} specifies the calendar month (1--12) on which the fiscal year
begins. The default is~1 (January), which makes fiscal and calendar years identical.

Fiscal quarter dates (\texttt{FYnn-Qk}) map to calendar dates as follows. \texttt{FY}$nn$
denotes the fiscal year beginning on month \texttt{fiscal\_year\_start} of calendar year
$20nn$. Fiscal quarter $k$ ($k \in \{1,2,3,4\}$) spans three consecutive calendar months
starting at calendar month $m_k$, where:
\[
  m_k = ((\texttt{fiscal\_year\_start} - 1) + (k - 1) \times 3) \bmod 12 + 1
\]
and the calendar year is $20nn$ if $m_k \ge \texttt{fiscal\_year\_start}$, otherwise
$20nn + 1$.

\textbf{Examples with \texttt{fiscal\_year\_start: 1} (default):}
FY26-Q1 = Jan--Mar 2026; FY26-Q2 = Apr--Jun 2026 (identical to calendar quarters).

\textbf{Examples with \texttt{fiscal\_year\_start: 4} (April fiscal year):}
FY26-Q1 = Apr--Jun 2026; FY26-Q2 = Jul--Sep 2026; FY26-Q3 = Oct--Dec 2026;
FY26-Q4 = Jan--Mar 2027.

If any \texttt{FY...} date appears in the document and \texttt{fiscal\_year\_start}
$\neq 1$, authors \textbf{should} set \texttt{fiscal\_year\_start} explicitly to ensure
all renderers produce consistent x-coordinates for fiscal dates.


%% ============================================================================
%% DATE MODEL
%% ============================================================================
\subsection{Date Model}
\label{sec:date-model}

Timelines present temporal information at varying granularities. A product roadmap may show
quarters; a conference schedule shows hours. The IR must represent this heterogeneity without
forcing false precision or losing meaningful vagueness.

\subsubsection{Date Representations}

A \texttt{date} in the IR is one of the following forms:

\begin{table}[h]
\centering
\small
\begin{tabular}{llp{6cm}}
\toprule
\textbf{Form} & \textbf{Example} & \textbf{Semantics} \\
\midrule
ISO date & \texttt{2026-06-09} & Specific calendar day \\
ISO datetime & \texttt{2026-06-09T14:00} & Specific moment (minute precision) \\
Year-month & \texttt{2026-06} & Entire month (June 2026) \\
Year only & \texttt{2026} & Entire year \\
Quarter & \texttt{2026-Q2} & Calendar quarter (Apr--Jun) \\
Half & \texttt{2026-H1} & Calendar half (Jan--Jun) \\
Fiscal quarter & \texttt{FY26-Q2} & Fiscal quarter (calendar mapping per \texttt{fiscal\_year\_start}; see \S\ref{sec:date-anchor}) \\
Relative & \texttt{+3m} & 3 months from the date anchor (\texttt{metadata.today} $\to$ \texttt{metadata.created}; see \S\ref{sec:date-anchor}) \\
Symbolic & \texttt{now} & Resolves to the date anchor (\texttt{metadata.today} $\to$ \texttt{metadata.created}; see \S\ref{sec:date-anchor}) \\
\bottomrule
\end{tabular}
\caption{Date representation forms}
\label{tab:date-forms}
\end{table}

\textbf{Design rationale:} We deliberately support coarse granularity (quarters, halves, years)
as first-class concepts rather than forcing conversion to day-level dates. This preserves
authorial intent: ``Q3 2026'' means the entire quarter, not ``2026-07-01''.

\subsubsection{Uncertain and Open Dates}

For dates that are genuinely unknown, tentative, or open-ended, the IR provides explicit
encodings rather than using null-as-magic:

\begin{table}[h]
\centering
\begin{tabular}{lp{8cm}}
\toprule
\textbf{Value} & \textbf{Semantics} \\
\midrule
\texttt{tbd} & Date is not yet determined; renders as ``TBD'' \\
\texttt{ongoing} & No planned end; activity continues indefinitely \\
\texttt{unknown} & Date exists but is not known to the author \\
\texttt{\textasciitilde 2026-Q3} & Approximate/tentative (prefix tilde) \\
\bottomrule
\end{tabular}
\caption{Uncertain and open date markers}
\label{tab:uncertain-dates}
\end{table}

\textbf{Invariant:} \texttt{tbd}, \texttt{ongoing}, and \texttt{unknown} are valid only in
positions where the schema allows uncertainty (e.g., \texttt{end} of a daterange, not
\texttt{start} of the document's \texttt{time\_range}).

\subsubsection{DateRange}

A \texttt{daterange} represents a temporal interval:

\begin{lstlisting}[language=yaml]
# Explicit start and end
start: 2026-04-01
end: 2026-06-30

# Open-ended (no planned end)
start: 2026-04-01
end: ongoing

# Shorthand for single-granularity span
span: 2026-Q2    # equivalent to start: 2026-04-01, end: 2026-06-30
\end{lstlisting}

\begin{table}[h]
\centering
\begin{tabular}{llccp{5cm}}
\toprule
\textbf{Field} & \textbf{Type} & \textbf{Req} & \textbf{Default} & \textbf{Description} \\
\midrule
\texttt{start} & \texttt{date} & Yes* & --- & Interval start (inclusive) \\
\texttt{end} & \texttt{date|ongoing|tbd} & No & \texttt{ongoing} & Interval end (inclusive); omitting is equivalent to \texttt{end: ongoing} \\
\texttt{span} & \texttt{date} & Yes* & --- & Single-unit shorthand; mutually exclusive with \texttt{start}/\texttt{end} \\
\bottomrule
\end{tabular}
\caption{DateRange fields (*exactly one of \texttt{span} or \texttt{start} (+ optional \texttt{end}) is required; co-presence of \texttt{span} with \texttt{start} or \texttt{end} is a well-formedness error)}
\label{tab:daterange-fields}
\end{table}

\textbf{Invariant:} When both \texttt{start} and \texttt{end} are concrete dates,
\texttt{end >= start}. The \texttt{span} shorthand expands to the natural bounds of its
granularity (e.g., \texttt{span: 2026-Q2} expands to Q2's first and last days).

\textbf{Invariant (exclusivity):} \texttt{span} is \emph{mutually exclusive} with
\texttt{start} and \texttt{end}. A daterange must use exactly one of:
\begin{itemize}
  \item \texttt{span} alone (shorthand form), or
  \item \texttt{start} alone or \texttt{start} + \texttt{end} (explicit form).
\end{itemize}
Co-presence of \texttt{span} with \texttt{start} or \texttt{end} on the same entity is
a well-formedness error (\texttt{SPAN\_START\_CONFLICT}).

\textbf{Open/ongoing semantics:} An Activity or daterange with \texttt{start} specified
but \texttt{end} omitted is an \emph{open/ongoing interval}. This is semantically
identical to writing \texttt{end: ongoing} and is an explicit, valid state (not an
authoring error). Renderers must extend the bar to the canvas right edge and apply
an open-end indicator (e.g., right-pointing chevron).


%% ============================================================================
%% TRACKS
%% ============================================================================
\subsection{Tracks (Swimlanes)}
\label{sec:ir-tracks}

Tracks are horizontal lanes that organize activities visually. Every activity belongs to
exactly one track. Tracks appear in document order (first track at top).

\begin{lstlisting}[language=yaml,caption={Track examples}]
tracks:
  - id: platform
    label: Platform Team
    description: Core infrastructure work

  - id: mobile
    label: Mobile Apps
    color: blue         # optional hint
\end{lstlisting}

\begin{table}[h]
\centering
\small
\begin{tabular}{llccp{4.5cm}}
\toprule
\textbf{Field} & \textbf{Type} & \textbf{Req} & \textbf{Default} & \textbf{Description} \\
\midrule
\texttt{id} & \texttt{id} & Yes & --- & Unique track identifier \\
\texttt{label} & \texttt{string} & Yes & --- & Display label \\
\texttt{description} & \texttt{string?} & No & \texttt{null} & Extended description \\
\texttt{color} & \texttt{string?} & No & theme & Color hint (theme may override) \\
\texttt{group} & \texttt{ref<Group>?} & No & \texttt{null} & Parent group (for nested tracks) \\
\texttt{index} & \texttt{int?} & No & position & Explicit sort index; unique non-negative integer; tracks render top-to-bottom in ascending \texttt{index} order \\
\texttt{collapsed} & \texttt{bool} & No & \texttt{false} & Initial collapsed state (if renderer supports) \\
\bottomrule
\end{tabular}
\caption{Track fields}
\label{tab:track-fields}
\end{table}

\textbf{Invariant:} Track \texttt{id} values must be unique within the document.

%% ============================================================================
%% GROUPS
%% ============================================================================
\subsection{Groups}
\label{sec:ir-groups}

Groups provide logical clustering of activities, tracks, or other groups. They enable
hierarchical organization (e.g., ``Engineering'' containing ``Platform'' and ``Mobile'' tracks)
without conflating logical grouping with visual swimlanes.

\begin{lstlisting}[language=yaml,caption={Group example}]
groups:
  - id: engineering
    label: Engineering
    children:
      - platform    # ref to track
      - mobile      # ref to track

  - id: releases
    label: Release Milestones
    members:
      - alpha-release   # ref to milestone
      - beta-release
\end{lstlisting}

\begin{table}[h]
\centering
\small
\begin{tabular}{llccp{4.5cm}}
\toprule
\textbf{Field} & \textbf{Type} & \textbf{Req} & \textbf{Default} & \textbf{Description} \\
\midrule
\texttt{id} & \texttt{id} & Yes & --- & Unique group identifier \\
\texttt{label} & \texttt{string} & Yes & --- & Display label \\
\texttt{description} & \texttt{string?} & No & \texttt{null} & Extended description \\
\texttt{parent} & \texttt{ref<Group>?} & No & \texttt{null} & Parent group (for nesting) \\
\texttt{children} & \texttt{list<ref<Track|Group>>} & No & \texttt{[]} & Ordered child tracks/groups \\
\texttt{members} & \texttt{list<ref<Activity|Milestone>>} & No & \texttt{[]} & Member entities \\
\texttt{color} & \texttt{string?} & No & theme & Color hint \\
\bottomrule
\end{tabular}
\caption{Group fields}
\label{tab:group-fields}
\end{table}

\textbf{Invariant:} Group hierarchies must be acyclic. A group cannot be its own ancestor.


%% ============================================================================
%% ACTIVITIES
%% ============================================================================
\subsection{Activities}
\label{sec:ir-activities}

Activities are time-spanning items---phases, projects, tasks, or efforts. They are the
primary content of most timelines.

\begin{lstlisting}[language=yaml,caption={Activity examples}]
activities:
  - id: api-redesign
    label: API Redesign
    track: platform
    start: 2026-Q1
    end: 2026-Q2
    status: in-progress
    progress: 0.4

  - id: mobile-launch
    label: Mobile App Launch
    track: mobile
    span: 2026-Q3
    status: planned
    category: launch
\end{lstlisting}

\begin{table}[h]
\centering
\small
\begin{tabular}{llccp{4cm}}
\toprule
\textbf{Field} & \textbf{Type} & \textbf{Req} & \textbf{Default} & \textbf{Description} \\
\midrule
\texttt{id} & \texttt{id} & Yes & --- & Unique activity identifier \\
\texttt{label} & \texttt{string} & Yes & --- & Display label (shown on bar) \\
\texttt{track} & \texttt{ref<Track>} & Yes & --- & Containing track \\
\texttt{start} & \texttt{date} & Yes* & --- & Start date \\
\texttt{end} & \texttt{date|ongoing|tbd} & No & \texttt{ongoing} & End date; omitting is equivalent to \texttt{end: ongoing} (open interval) \\
\texttt{span} & \texttt{date} & Yes* & --- & Single-unit shorthand; mutually exclusive with \texttt{start}/\texttt{end} \\
\texttt{status} & \texttt{Status} & No & \texttt{planned} & Visual communication status \\
\texttt{progress} & \texttt{number[0..1]?} & No & \texttt{null} & Completion fraction \\
\texttt{category} & \texttt{string?} & No & \texttt{null} & Semantic category (for styling) \\
\texttt{description} & \texttt{string?} & No & \texttt{null} & Extended description \\
\texttt{group} & \texttt{ref<Group>?} & No & \texttt{null} & Logical group membership \\
\texttt{url} & \texttt{string?} & No & \texttt{null} & External link \\
\texttt{tags} & \texttt{list<string>} & No & \texttt{[]} & Freeform tags \\
\texttt{metadata} & \texttt{map<string,any>} & No & \texttt{\{\}} & Extension metadata \\
\bottomrule
\end{tabular}
\caption{Activity fields (*exactly one of \texttt{span} or \texttt{start} (+ optional \texttt{end}) is required; co-presence of \texttt{span} with \texttt{start} or \texttt{end} is a well-formedness error)}
\label{tab:activity-fields}
\end{table}

\subsubsection{Status Enum}
\label{sec:status-enum}

The \texttt{status} field communicates \emph{visual} state for presentation purposes.
This is explicitly \textbf{not} a project-management workflow state; it describes what
the audience should understand about the item's current condition.

\begin{lstlisting}
enum Status {
  planned,      # Scheduled but not yet started
  in-progress,  # Currently active
  done,         # Completed
  at-risk,      # May miss target (visual warning)
  blocked,      # Cannot proceed (visual alert)
  cancelled,    # No longer planned
  tentative     # Not yet confirmed
}
\end{lstlisting}

\textbf{Design rationale:} We include \texttt{at-risk} and \texttt{blocked} because executive
timelines frequently need to communicate risk visually. We omit workflow states like
``in review'' or ``approved'' which belong to project-management tools, not visual
communication artifacts.

\subsubsection{Progress}

The \texttt{progress} field is a number in $[0, 1]$ representing completion fraction.
It has no inherent visual meaning---the theme decides whether to render it as a
progress bar, percentage label, or ignore it entirely.

\textbf{Invariant:} If \texttt{progress} is provided, it must be in $[0, 1]$.
\texttt{progress: 1.0} does not imply \texttt{status: done}; these are orthogonal.


%% ============================================================================
%% MILESTONES
%% ============================================================================
\subsection{Milestones}
\label{sec:ir-milestones}

Milestones are point-in-time markers---releases, deadlines, decision points, or events.
Unlike activities, they have no duration.

\begin{lstlisting}[language=yaml,caption={Milestone examples}]
milestones:
  - id: alpha-release
    label: Alpha Release
    date: 2026-03-15
    track: platform
    status: done
    icon: flag

  - id: board-review
    label: Board Review
    date: 2026-Q2        # first day of Q2
    category: governance
\end{lstlisting}

\begin{table}[h]
\centering
\small
\begin{tabular}{llccp{4.5cm}}
\toprule
\textbf{Field} & \textbf{Type} & \textbf{Req} & \textbf{Default} & \textbf{Description} \\
\midrule
\texttt{id} & \texttt{id} & Yes & --- & Unique milestone identifier \\
\texttt{label} & \texttt{string} & Yes & --- & Display label \\
\texttt{date} & \texttt{date} & Yes & --- & Point in time \\
\texttt{track} & \texttt{ref<Track>?} & No & global & Associated track (or global) \\
\texttt{status} & \texttt{Status} & No & \texttt{planned} & Visual status \\
\texttt{category} & \texttt{string?} & No & \texttt{null} & Semantic category \\
\texttt{icon} & \texttt{string?} & No & theme & Icon hint (diamond, flag, star, etc.) \\
\texttt{description} & \texttt{string?} & No & \texttt{null} & Extended description \\
\texttt{group} & \texttt{ref<Group>?} & No & \texttt{null} & Logical group membership \\
\texttt{url} & \texttt{string?} & No & \texttt{null} & External link \\
\texttt{tags} & \texttt{list<string>} & No & \texttt{[]} & Freeform tags \\
\bottomrule
\end{tabular}
\caption{Milestone fields}
\label{tab:milestone-fields}
\end{table}

\textbf{Note:} When \texttt{track} is omitted, the milestone is ``global'' and may be
rendered spanning all tracks (e.g., a vertical line with label at top).

%% ============================================================================
%% ANNOTATIONS
%% ============================================================================
\subsection{Annotations}
\label{sec:ir-annotations}

Annotations add contextual information---callouts, notes, highlights, or today-markers---
without being first-class timeline entities.

\begin{lstlisting}[language=yaml,caption={Annotation examples}]
annotations:
  - type: today-marker
    date: now

  - type: callout
    target: api-redesign
    text: "Critical path item"
    position: above

  - type: bracket
    targets: [alpha-release, beta-release]
    label: "Testing Phase"

  - type: period
    start: 2026-06-01
    end: 2026-06-15
    label: "Code Freeze"
    style: highlight
\end{lstlisting}

\begin{table}[h]
\centering
\small
\begin{tabular}{llccp{4cm}}
\toprule
\textbf{Field} & \textbf{Type} & \textbf{Req} & \textbf{Default} & \textbf{Description} \\
\midrule
\texttt{type} & \texttt{AnnotationType} & Yes & --- & Annotation kind \\
\texttt{id} & \texttt{id?} & No & auto & Optional identifier \\
\texttt{target} & \texttt{ref<Activity|Milestone>?} & Cond & --- & Single target entity \\
\texttt{targets} & \texttt{list<ref<...>>?} & Cond & --- & Multiple targets \\
\texttt{date} & \texttt{date?} & Cond & --- & Point annotation \\
\texttt{start}/\texttt{end} & \texttt{date?} & Cond & --- & Range annotation \\
\texttt{text}/\texttt{label} & \texttt{string?} & No & \texttt{null} & Display text \\
\texttt{position} & \texttt{enum\{above,below,left,right\}?} & No & auto & Hint only \\
\texttt{style} & \texttt{string?} & No & theme & Style hint \\
\bottomrule
\end{tabular}
\caption{Annotation fields (conditional fields depend on type)}
\label{tab:annotation-fields}
\end{table}

\subsubsection{Annotation Types}

\begin{lstlisting}
enum AnnotationType {
  today-marker,  # Vertical line at current date
  callout,       # Note attached to an entity
  bracket,       # Visual bracket spanning entities
  period,        # Highlighted time period
  note,          # Freestanding note
  connector      # Line connecting two entities
}
\end{lstlisting}


%% ============================================================================
%% SECTIONS
%% ============================================================================
\subsection{Sections}
\label{sec:ir-sections}

Sections divide the timeline into logical or visual chapters. They enable multi-page
timelines or distinct phases that should be visually separated.

\begin{lstlisting}[language=yaml,caption={Section example}]
sections:
  - id: current-state
    label: "Current State (Q1-Q2)"
    time_range:
      start: 2026-Q1
      end: 2026-Q2

  - id: future-state
    label: "Future Roadmap (Q3-Q4)"
    time_range:
      start: 2026-Q3
      end: 2026-Q4
\end{lstlisting}

\begin{table}[h]
\centering
\small
\begin{tabular}{llccp{5cm}}
\toprule
\textbf{Field} & \textbf{Type} & \textbf{Req} & \textbf{Default} & \textbf{Description} \\
\midrule
\texttt{id} & \texttt{id} & Yes & --- & Unique section identifier \\
\texttt{label} & \texttt{string} & Yes & --- & Section title \\
\texttt{time\_range} & \texttt{TimeRange?} & No & inherit & Override document time range \\
\texttt{tracks} & \texttt{list<ref<Track>>?} & No & all & Subset of tracks to show \\
\texttt{description} & \texttt{string?} & No & \texttt{null} & Section description \\
\texttt{page\_break} & \texttt{bool} & No & \texttt{false} & Hint for paginated output \\
\bottomrule
\end{tabular}
\caption{Section fields}
\label{tab:section-fields}
\end{table}

%% ============================================================================
%% LEGEND
%% ============================================================================
\subsection{Legend}
\label{sec:ir-legend}

The legend documents the meaning of statuses, categories, and visual conventions used
in the timeline. If omitted, the renderer may auto-generate a legend from used values.

\begin{lstlisting}[language=yaml,caption={Legend example}]
legend:
  show: true
  position: bottom-right
  entries:
    - key: in-progress
      label: In Progress
      description: Currently being worked on
    - key: at-risk
      label: At Risk
      description: May miss planned date
    - key: launch
      label: Launch Event
      category: true
\end{lstlisting}

\begin{table}[h]
\centering
\small
\begin{tabular}{llccp{5cm}}
\toprule
\textbf{Field} & \textbf{Type} & \textbf{Req} & \textbf{Default} & \textbf{Description} \\
\midrule
\texttt{show} & \texttt{bool} & No & \texttt{true} & Whether to render legend \\
\texttt{position} & \texttt{string?} & No & theme & Position hint \\
\texttt{title} & \texttt{string?} & No & \texttt{"Legend"} & Legend title \\
\texttt{entries} & \texttt{list<LegendEntry>} & No & auto & Legend entries \\
\bottomrule
\end{tabular}
\caption{Legend fields}
\label{tab:legend-fields}
\end{table}

\begin{table}[h]
\centering
\small
\begin{tabular}{llccp{5cm}}
\toprule
\textbf{Field} & \textbf{Type} & \textbf{Req} & \textbf{Default} & \textbf{Description} \\
\midrule
\texttt{key} & \texttt{string} & Yes & --- & Status or category key \\
\texttt{label} & \texttt{string} & Yes & --- & Human-readable label \\
\texttt{description} & \texttt{string?} & No & \texttt{null} & Extended description \\
\texttt{category} & \texttt{bool} & No & \texttt{false} & True if this is a category (not status) \\
\bottomrule
\end{tabular}
\caption{LegendEntry fields}
\label{tab:legend-entry-fields}
\end{table}


%% ============================================================================
%% ID AND REFERENCE SYSTEM
%% ============================================================================
\subsection{ID and Reference System}
\label{sec:ir-refs}

\subsubsection{Identifier Format}

All \texttt{id} fields must be:
\begin{itemize}
  \item Non-empty strings
  \item Composed of lowercase letters, digits, and hyphens
  \item Starting with a letter
  \item Unique within their entity type AND globally within the document
\end{itemize}

\textbf{Regex:} \verb|^[a-z][a-z0-9-]*$|

\textbf{Rationale for global uniqueness:} References may appear in contexts where the
entity type is ambiguous (e.g., annotation targets). Global uniqueness eliminates
the need for type prefixes in references.

\begin{lstlisting}[language=yaml,caption={ID examples}]
# Valid IDs
id: api-v2
id: q3-release
id: mobile-team-ios

# Invalid IDs
id: API-v2         # uppercase
id: 2026-release   # starts with digit
id: api_v2         # underscore
\end{lstlisting}

\subsubsection{References}

A reference is a string containing the \texttt{id} of another entity. The IR uses
bare strings for references (not structured objects) for YAML terseness:

\begin{lstlisting}[language=yaml]
activities:
  - id: feature-a
    track: platform      # ref<Track>
    group: engineering   # ref<Group>

annotations:
  - type: callout
    target: feature-a    # ref<Activity>
\end{lstlisting}

\textbf{Invariant:} Every reference must resolve to exactly one entity in the document.
A validator must reject documents with dangling references.

\subsubsection{Reference Namespacing}

Because IDs are globally unique, references are unambiguous without type prefixes.
The schema context determines what types are valid:

\begin{itemize}
  \item \texttt{activity.track}: must reference a Track
  \item \texttt{annotation.target}: may reference Activity or Milestone
  \item \texttt{group.children}: may reference Track or Group
\end{itemize}

A validator checks that the referenced entity has the correct type for its context.

%% ============================================================================
%% WELL-FORMEDNESS AND INVARIANTS
%% ============================================================================
\subsection{Well-Formedness Rules}
\label{sec:ir-invariants}

A Timeline IR document is \textbf{well-formed} if and only if all the following
invariants hold:

\subsubsection{Structural Invariants}

\begin{enumerate}
  \item \textbf{Version present:} \texttt{version} field exists and matches a known
        specification version.
  \item \textbf{Required fields:} All fields marked ``Req: Yes'' are present and non-null.
  \item \textbf{At least one track:} The \texttt{tracks} list is non-empty.
  \item \textbf{Unique IDs:} All \texttt{id} values are globally unique within the document.
  \item \textbf{Valid ID format:} All \texttt{id} values match \verb|^[a-z][a-z0-9-]*$|.
\end{enumerate}

\subsubsection{Referential Invariants}

\begin{enumerate}[resume]
  \item \textbf{References resolve:} Every reference points to an existing entity.
  \item \textbf{Reference types match:} References point to entities of the expected type
        (e.g., \texttt{activity.track} references a Track, not a Group).
  \item \textbf{No circular groups:} The group parent-child graph is acyclic.
\end{enumerate}

\subsubsection{Temporal Invariants}

\begin{enumerate}[resume]
  \item \textbf{Valid time range:} If \texttt{metadata.time\_range} has both start and end,
        then \texttt{end >= start}.
  \item \textbf{Activity duration non-negative:} For activities with concrete start and end,
        \texttt{end >= start}.
  \item \textbf{Date format valid:} All date values conform to the date model
        (\S\ref{sec:date-model}).
  \item \textbf{Activity date-source exclusive:} Every activity must satisfy exactly one of:
        (a)~\texttt{span} is present and neither \texttt{start} nor \texttt{end} is present; or
        (b)~\texttt{start} is present and \texttt{span} is absent (\texttt{end} optional).
        Co-presence of \texttt{span} with \texttt{start} or \texttt{end} is a hard
        well-formedness error: \texttt{SPAN\_START\_CONFLICT}.
  \item \textbf{Date anchor present when required:} If any date field in the document
        contains \texttt{now} or a relative offset (e.g., \texttt{+3m}, \texttt{-2w}),
        then at least one of \texttt{metadata.today} or \texttt{metadata.created} must be
        present. Absence of both is a hard error: \texttt{DATE\_ANCHOR\_MISSING}.
  \item \textbf{Track index values unique:} When \texttt{track.index} is present, all
        supplied values must be unique non-negative integers within the document.
\end{enumerate}

\subsubsection{Value Invariants}

\begin{enumerate}[resume]
  \item \textbf{Progress in range:} If \texttt{progress} is present, $0 \le \text{progress} \le 1$.
  \item \textbf{Status valid:} All \texttt{status} values are members of the Status enum.
  \item \textbf{Axis unit valid:} If present, \texttt{axis\_unit} is a member of AxisUnit enum.
\end{enumerate}

\subsubsection{Soft Warnings (Valid but Suspicious)}

The following are not errors but may indicate authorial mistakes:

\begin{itemize}
  \item Activity or milestone outside \texttt{metadata.time\_range} (will be clipped)
  \item \texttt{progress: 1.0} with \texttt{status: in-progress} (likely stale)
  \item Empty \texttt{activities} and \texttt{milestones} lists (vacuous timeline)
  \item Unused tracks (no activities reference them)
\end{itemize}


%% ============================================================================
%% EXAMPLES
%% ============================================================================
\subsection{Complete Examples}
\label{sec:ir-examples}

The following examples demonstrate realistic Timeline IR documents. Each is a complete,
valid document showcasing different timeline use cases.

\subsubsection{Example 1: Product Roadmap}

A quarterly product roadmap with multiple teams, showing how the IR represents a
typical software planning artifact.

\begin{lstlisting}[language=yaml,caption={Product roadmap example},label={lst:roadmap-example}]
version: "1.0"

metadata:
  title: "Product Roadmap 2026"
  subtitle: "Engineering & Design"
  author: "Product Team"
  created: 2026-06-09
  time_range:
    start: 2026-Q1
    end: 2026-Q4
  axis_unit: quarter
  theme: corporate-blue

tracks:
  - id: platform
    label: Platform
    description: Core infrastructure and APIs

  - id: mobile
    label: Mobile Apps
    description: iOS and Android applications

  - id: design
    label: Design System
    description: Component library and guidelines

groups:
  - id: foundation
    label: Foundation Work
    members: [api-v2, component-lib]

activities:
  - id: api-v2
    label: API v2 Migration
    track: platform
    start: 2026-Q1
    end: 2026-Q2
    status: in-progress
    progress: 0.6
    category: infrastructure

  - id: mobile-redesign
    label: Mobile App Redesign
    track: mobile
    start: 2026-Q2
    end: 2026-Q3
    status: planned
    description: Complete visual refresh

  - id: component-lib
    label: Component Library
    track: design
    start: 2026-Q1
    end: 2026-Q4
    status: in-progress
    progress: 0.3

  - id: analytics
    label: Analytics Integration
    track: platform
    span: 2026-Q3
    status: planned

milestones:
  - id: api-ga
    label: API v2 GA
    date: 2026-06-30
    track: platform
    status: planned
    icon: flag

  - id: app-launch
    label: App Store Launch
    date: 2026-Q4
    track: mobile
    status: planned
    category: launch

annotations:
  - type: today-marker
    date: now

legend:
  show: true
  entries:
    - key: infrastructure
      label: Infrastructure
      category: true
    - key: launch
      label: Launch Event
      category: true
\end{lstlisting}

\subsubsection{Example 2: Executive Program Timeline}

A high-level transformation program timeline suitable for board presentations,
demonstrating status communication and risk visibility.

\begin{lstlisting}[language=yaml,caption={Executive program timeline},label={lst:exec-timeline}]
version: "1.0"

metadata:
  title: "Digital Transformation Program"
  subtitle: "FY26 Executive View"
  author: "PMO"
  created: 2026-06-09
  time_range:
    start: 2026-01-01
    end: 2026-12-31
  axis_unit: month
  theme: executive

tracks:
  - id: strategy
    label: Strategy & Governance

  - id: technology
    label: Technology Modernization

  - id: people
    label: People & Change

activities:
  - id: strategy-define
    label: Strategy Definition
    track: strategy
    start: 2026-01-01
    end: 2026-03-31
    status: done
    progress: 1.0

  - id: vendor-selection
    label: Vendor Selection
    track: technology
    start: 2026-02-01
    end: 2026-04-30
    status: done

  - id: platform-build
    label: Platform Build
    track: technology
    start: 2026-04-01
    end: 2026-09-30
    status: in-progress
    progress: 0.35

  - id: data-migration
    label: Data Migration
    track: technology
    start: 2026-07-01
    end: 2026-11-30
    status: at-risk
    description: Dependency on legacy system documentation

  - id: training
    label: Training Program
    track: people
    start: 2026-06-01
    end: 2026-12-31
    status: planned

  - id: change-mgmt
    label: Change Management
    track: people
    start: 2026-03-01
    end: ongoing
    status: in-progress

milestones:
  - id: board-approval
    label: Board Approval
    date: 2026-03-15
    status: done
    icon: star

  - id: go-live
    label: Go Live
    date: 2026-10-01
    track: technology
    status: planned
    icon: flag

  - id: program-close
    label: Program Close
    date: 2026-12-15
    status: planned

annotations:
  - type: today-marker
    date: now

  - type: callout
    target: data-migration
    text: "Risk: Legacy documentation gaps"
    position: above

  - type: bracket
    targets: [platform-build, data-migration]
    label: "Critical Path"

sections:
  - id: h1
    label: "H1 2026: Foundation"
    time_range:
      start: 2026-01-01
      end: 2026-06-30

  - id: h2
    label: "H2 2026: Execution"
    time_range:
      start: 2026-07-01
      end: 2026-12-31
\end{lstlisting}


\subsubsection{Example 3: Architecture Evolution Timeline}

A technical architecture timeline showing system evolution over multiple years,
demonstrating coarse granularity and technical categorization.

\begin{lstlisting}[language=yaml,caption={Architecture evolution timeline},label={lst:arch-timeline}]
version: "1.0"

metadata:
  title: "Platform Architecture Evolution"
  subtitle: "2026-2028 Technical Roadmap"
  author: "Architecture Team"
  created: 2026-06-09
  time_range:
    start: 2026
    end: 2028
  axis_unit: half
  theme: technical

tracks:
  - id: frontend
    label: Frontend

  - id: backend
    label: Backend Services

  - id: data
    label: Data Platform

  - id: infra
    label: Infrastructure

activities:
  - id: react-migration
    label: React 19 Migration
    track: frontend
    start: 2026-H1
    end: 2026-H2
    status: in-progress
    category: modernization

  - id: micro-frontend
    label: Micro-Frontend Architecture
    track: frontend
    start: 2027-H1
    end: 2027-H2
    status: planned
    category: architecture

  - id: api-gateway
    label: API Gateway
    track: backend
    span: 2026-H1
    status: done
    category: infrastructure

  - id: service-mesh
    label: Service Mesh Adoption
    track: backend
    start: 2026-H2
    end: 2027-H1
    status: planned
    category: infrastructure

  - id: event-driven
    label: Event-Driven Architecture
    track: backend
    start: 2027
    end: 2028
    status: tentative
    category: architecture

  - id: lakehouse
    label: Data Lakehouse
    track: data
    start: 2026-H1
    end: 2027-H1
    status: in-progress
    progress: 0.4
    category: data

  - id: ml-platform
    label: ML Platform
    track: data
    start: 2027
    end: ongoing
    status: tentative
    category: ai

  - id: k8s-migration
    label: Kubernetes Migration
    track: infra
    span: 2026
    status: in-progress
    progress: 0.7
    category: infrastructure

  - id: multi-cloud
    label: Multi-Cloud Strategy
    track: infra
    start: 2027
    end: 2028
    status: planned
    category: infrastructure

milestones:
  - id: legacy-sunset
    label: Legacy System Sunset
    date: 2027-H1
    track: backend
    status: planned
    icon: flag

  - id: cloud-native
    label: Cloud Native Milestone
    date: 2026
    track: infra
    status: planned

annotations:
  - type: period
    start: 2026-H2
    end: 2027-H1
    label: "Transition Period"
    style: highlight

legend:
  entries:
    - key: modernization
      label: Modernization
      category: true
    - key: architecture
      label: Architecture Change
      category: true
    - key: infrastructure
      label: Infrastructure
      category: true
    - key: data
      label: Data Platform
      category: true
    - key: ai
      label: AI/ML
      category: true
\end{lstlisting}


%% ============================================================================
%% EXTENSIBILITY
%% ============================================================================
\subsection{Extensibility}
\label{sec:ir-extensibility}

The IR is designed to be extended without breaking existing consumers.

\subsubsection{Metadata Extension}

Every entity type includes a \texttt{metadata} field (type \texttt{map<string,any>})
for application-specific data:

\begin{lstlisting}[language=yaml]
activities:
  - id: feature-x
    label: Feature X
    track: platform
    span: 2026-Q2
    metadata:
      jira_key: PROJ-1234
      owner: alice@example.com
      priority: high
\end{lstlisting}

Renderers should ignore unrecognized metadata keys. Source integrations may use
metadata to preserve round-trip fidelity with external systems.

\subsubsection{Custom Categories}

The \texttt{category} field accepts arbitrary strings. The theme engine maps
categories to visual styles. New categories can be introduced without IR changes:

\begin{lstlisting}[language=yaml]
activities:
  - id: security-audit
    category: security   # custom category
    # ...

legend:
  entries:
    - key: security
      label: Security Work
      category: true
\end{lstlisting}

\subsubsection{Version Compatibility}

The \texttt{version} field enables forward compatibility:

\begin{itemize}
  \item Consumers should accept documents with higher minor versions (1.1, 1.2, etc.)
        and ignore unknown fields.
  \item Major version changes (2.0) may introduce breaking changes; consumers should
        reject documents with unsupported major versions.
\end{itemize}

%% ============================================================================
%% SERIALIZATION
%% ============================================================================
\subsection{Serialization and Syntax}
\label{sec:ir-serialization}

\subsubsection{Canonical Format}

YAML 1.2 is the canonical surface syntax for Timeline IR documents. Design decisions
prioritizing YAML:

\begin{itemize}
  \item \textbf{Minimal punctuation:} No braces for simple mappings, no quotes for
        most strings.
  \item \textbf{Line-oriented:} Each field on its own line for clean diffs.
  \item \textbf{Stable ordering:} Fields should appear in specification order for
        consistency; validators may enforce this.
  \item \textbf{Comments preserved:} YAML allows comments; tools should preserve them.
\end{itemize}

\subsubsection{Alternative Formats}

Tooling may support additional formats with lossless conversion:

\begin{itemize}
  \item \textbf{JSON:} Direct mapping; verbose but universally parseable.
  \item \textbf{TOML:} Alternative for users who prefer it.
\end{itemize}

The IR is defined abstractly; the YAML syntax is one serialization, not the definition.

\subsubsection{Git Friendliness}

For minimal diff churn:

\begin{enumerate}
  \item Lists use block style (one item per line), not flow style.
  \item Fields appear in consistent order.
  \item Auto-generated fields (like \texttt{updated} timestamps) are optional and
        may be omitted in version-controlled documents.
  \item IDs are stable slugs, not generated UUIDs.
\end{enumerate}

\begin{lstlisting}[language=yaml,caption={Git-friendly style}]
# Good: block style, stable order
activities:
  - id: feature-a
    label: Feature A
    track: platform
    span: 2026-Q2

# Avoid: flow style, hard to diff
activities: [{id: feature-a, label: Feature A, track: platform, span: 2026-Q2}]
\end{lstlisting}

%% ============================================================================
%% VALIDATION
%% ============================================================================
\subsection{Validation}
\label{sec:ir-validation}

A conforming validator must check all invariants in \S\ref{sec:ir-invariants}.
Validation is a function:

\[
\text{validate} : \text{Document} \to \text{Valid} \mid \text{Errors}
\]

\subsubsection{Error Reporting}

Validation errors should include:
\begin{itemize}
  \item Path to the offending field (e.g., \texttt{activities[2].track})
  \item The invariant violated
  \item The actual value
  \item Suggested fix (when determinable)
\end{itemize}

\subsubsection{Strict vs. Lenient Mode}

\begin{itemize}
  \item \textbf{Strict mode:} All invariants enforced; required for rendering.
  \item \textbf{Lenient mode:} Warnings instead of errors for soft issues; useful
        for work-in-progress documents.
\end{itemize}

%% ============================================================================
%% RELATIONSHIP TO OTHER COMPONENTS
%% ============================================================================
\subsection{Relationship to Other Components}
\label{sec:ir-relationships}

The IR is the stable contract between pipeline stages:

\begin{description}
  \item[Ingestion (upstream)] Source adapters (ADO, GitHub, Markdown, etc.) produce
        IR documents. The IR defines the target shape; adapters handle mapping.

  \item[Theme Engine (downstream)] Themes consume IR and produce styled visual
        specifications. Themes interpret \texttt{status}, \texttt{category}, and
        hint fields (\texttt{color}, \texttt{icon}) but cannot require additional
        IR fields.

  \item[Renderer (downstream)] The renderer consumes themed IR (or IR + theme) and
        produces output (SVG, PNG, PDF, PPTX). The IR remains rendering-agnostic;
        pixel-level decisions are not IR concerns.

  \item[Agent Integration] AI agents generate IR directly. The IR's explicit typing,
        clear required/optional fields, and forgiving structure enable reliable
        generation. Agents need not understand rendering.
\end{description}

\textbf{Contract stability:} Downstream components may not require IR fields beyond
this specification. Extension fields (\texttt{metadata}) are pass-through.

